\scanpage{40}
\begin{center}6 - 26 -\hfill n\textsuperscript{o} 392\end{center}

à l'intérieur de chaque intervalle, prendre $b_{n+1}=b_nz^{i-1}x^{-1}$ ; à l'extrémité de chaque intervalle, utiliser l'inégalité $ijk-ij-jk-ki\geq 0$).

b) On suppose que $z^{i-1}\notin Ax+Ay$. Montrer qu'il n'existe dans l'anneau $A[[T]]$ aucune série formelle de terme constant $y^k$, ($k$ entier $>0$), et qui soit un élément associé à $v$ dans $B[[T]]$, (calculer le coefficient de $T$ dans le produit de $v$ et d'un élément inversible de $B[[T]]$).

c) Déduire de a), b) et de l'exerc. 7 un exemple d'anneau factoriel $A$ tel que l'anneau de séries formelles $A[[T]]$ ne soit pas factoriel. (Avec les notations précédentes, et en supposant remplies les conditions de a) et b) sur $x$, $y$, $z$, $i$, $j$, $k$, montrer que $vv'$ ne peut être produit de facteurs extrémaux $u_h$, ($1\leq h\leq r$), dans $A[[T]]$ ; on observera que les $u_h$ sont des séries formelles dont les termes constants sont des puissances de $y$. Considérer ensuite l'anneau $C=S'^{-1}A[[T]]$, où $S'$ est formé des séries dont le terme constant est dans $S$ ; $B[[T]]$ est le complété de l'anneau de Zariski $C$ ; montrer que $v'\in C$, et que $v$ et les $u_h$ sont extrémaux dans $C$ ; obtenir finalement une contradiction avec b)).

d) Déduire de c) un exemple d'anneau local factoriel noethérien $A$ dont le complété $\widehat{A}$ est un anneau de Krull non factoriel.

9)\quad a) Soit $A$ un anneau intègre noethérien tel que pour tout idéal maximal $\underline{m}$ de $A$, $A_{\underline{m}}$ soit un anneau de valuation discrète. Si $B$ est l'anneau de séries formelles $A[[X_1,\ldots,X_n]]$, montrer que pour tout idéal maximal $\underline{n}$ de $B$, l'anneau $B_{\underline{n}}$ est factoriel, (considérer son complété, en utilisant la prop. 8 du n\textsuperscript{o} 9). En déduire que tout idéal divisoriel de $B$ est un $B$-module projectif.

b) Soit $C$ un anneau noethérien tel que tout $C$-module projectif de type fini soit libre ; montrer que l'anneau de séries formelles $C[[X]]$ a la même propriété (cf. chap. II, \S\ 3, n\textsuperscript{o} 2, prop. 5).

c) Déduire de a) et b) que si $A$ est un anneau principal, l'anneau de séries formelles $A[[X_1,\ldots,X_n]]$ est factoriel.

10)\quad a) Un anneau factoriel prüférien, (\S\ 2, exerc. 12), est principal.

b) Un anneau de Krull pseudo-bezoutien, (\S\ 1, exerc. 21), est factoriel.

11) Soient $K$ un corps, $A$ l'anneau de polynômes $K[X,Y]$, qui est factoriel ; si $L$ est le corps $K(X^2,Y/X)\subset K(X,Y)$, montrer que l'anneau $A\cap L$ n'est pas factoriel.

\margin{39}

